% 本章按原 Word 结构转换。\templatehint 为写作提示，可用文档类 final 选项隐藏。

\chapter{\texorpdfstring{绪\thesistitlespace 论}{绪论}}\label{chap:1}

\templatehint{{\centering\heiti\fontsize{15.96bp}{31.2bp}\selectfont\ecjtuwordstrut{21.4}{9.8}（以上为一级标题，黑体三号加粗，居中）\par}}
\vspace{19.44bp}
\templatehint{({\fontspec{Times New Roman}\XeTeXcharclass"2026=0\relax ……}..\hspace{0bp}此处为论文内容)}
\vspace{19.5bp}

\section{研究的背景及意义}

\templatehint{{\heiti\fontsize{14.04bp}{31.2bp}\selectfont\ecjtuwordstrut{20.68}{10.52}（以上为二级标题，黑体四号加粗，居左顶格）}}
\templatehint{({\fontspec{Times New Roman}\XeTeXcharclass"2026=0\relax ……}..\hspace{6bp}此处为论文内容)}

\subsection{选题的背景}

\templatehint{（以上为三级标题，黑体小四，居左顶格）}

\ifecjtuguide\vspace{0.08bp}\fi
随着计算机技术特别是互联网技术的不断发展，信息技术对人类社会生活的影响越来越大，正日益改变着我们的生活方式。作为社会经济的主体细胞——企业，无疑也在接受最为深刻的变革，从上个世纪\hspace{5.904bp}50\hspace{5.88bp}年代以来，信息技术不断进入到企业的日常运营管理中来，随着技术的进步，也逐渐深入到各种不同的企业管理应用中，慢慢形成了各种各样的企业应用系统，如TPS，MIS，MRP，ERP，\hspace{6bp}OA等。计算机技术本身也在经历着深刻的\linebreak 变化，无论是硬件还是软件层面都发生了巨大的变革，从硬件功能的不断增强，以及各种\linebreak 标准和规范的不断完善和统一，计算机技术不断向全球互联互通信息共享的方向发展。


企业管理是企业管理者为了实现某种预期目标，应用管理的思想理论和方法，合理地计划、组织、指挥、协调和控制人力资源，调度企业的各种资源，以求用最小的投入获得最大的产出。


{\xeCJKsetup{CJKglue=\hskip0bp}对人力资源进行合理配置、\hspace{-2.16bp}有效开发和科学管理，\hspace{-2.16bp}从而影响雇员的态度、\hspace{-2.16bp}行为和绩效，\\借以实现组织目标的管理活动。\par}

\templatehint{\hspace*{24bp}{\color{blue}{\fontspec{Times New Roman}\XeTeXcharclass"2026=0\relax ……}}（以上为范文示例）}


\subsection{国内外研究现状}

\templatehint{({\fontspec{Times New Roman}\XeTeXcharclass"2026=0\relax ……}..\hspace{6bp}此处为论文内容)}

% 学院原稿在此处分为绪论的第二页。
\clearpage
\subsection{研究的意义}

\ifecjtuguide\vspace{58.5bp}\fi

\section{系统目标}

\ifecjtuguide\vspace{39bp}\fi

